\documentclass[preview]{standalone}
\usepackage{amsmath}
\usepackage{amssymb}
\usepackage{stellar}
\usepackage{definitions}
\usepackage{bettelini}
\begin{document}
\id{diffeq-riccati}
\genpage
\section{Riccati equations}

\begin{snippetdefinition}{diffeq-riccati-definition}{Riccati equation}
    Let \(I\) be an open interval and \(a,b,c\in C(I)\).
    A \emph{Riccati equation} has the form
    \[
        y'+a(x)y+b(x)y^2+c(x)=0.
    \]
    If \(b\equiv0\), it reduces to a linear equation.
\end{snippetdefinition}

\begin{snippettheorem}{diffeq-riccati-reduction}{Reduction using a particular solution}
    Let \(a,b,c\in C(I)\), and let \(y_*\) solve
    \(y'+ay+by^2+c=0\) on an open interval \(I\).
    On any subinterval where \(y\neq y_*\), the substitution
    \[
        y=y_*+\frac1z
    \]
    transforms the equation into
    \[
        z'-\bigl(a+2by_*\bigr)z=b.
    \]
    Conversely, any nonvanishing solution \(z\) of this linear equation
    gives a Riccati solution by that formula.
    The solution \(y=y_*\) must be included separately.
\end{snippettheorem}

\begin{snippetproof}{diffeq-riccati-reduction-proof}{diffeq-riccati-reduction}{Reduction using a particular solution}
    Set \(y=y_*+u\). Expanding gives
    \begin{align*}
        0&=y_*'+u'+a(y_*+u)+b(y_*^2+2y_*u+u^2)+c\\
         &=\underbrace{y_*'+ay_*+by_*^2+c}_{0}
           +u'+(a+2by_*)u+bu^2.
    \end{align*}
    Thus \(u\) satisfies a Bernoulli equation of exponent two.
    Write \(A=a+2by_*\) and substitute \(u=1/z\):
    \[
        -\frac{z'}{z^2}+\frac A z+\frac b{z^2}=0
        \quad\Longleftrightarrow\quad z'-Az=b.
    \]
    All steps are reversible when \(z\neq0\).
    Moreover \(f=-ay-by^2-c\) and \(\partial_yf=-a-2by\) are continuous,
    so uniqueness implies that a solution differing initially from \(y_*\)
    cannot meet it later on their common interval.
\end{snippetproof}

\section{Examples}

\begin{snippetexample}{diffeq-riccati-cosine-example}{A constant particular solution}
    Consider \(y'+(y^2+y-\tfrac34)\cos x=0\), \(y(0)=y_0\).
    Since \((1/2)^2+1/2-3/4=0\), \(y_*=1/2\) is a particular solution.
    If \(y_0=1/2\), it is the required global solution.
    Otherwise set \(y=1/2+1/z\), obtaining
    \[
        z'-2\cos x\,z=\cos x,\qquad z(0)=\frac1{y_0-1/2}.
    \]
    Multiplying by \(e^{-2\sin x}\) gives
    \[
        (ze^{-2\sin x})'=\cos x\,e^{-2\sin x},
        \qquad
        z(x)=\left(\frac1{y_0-1/2}+\frac12\right)e^{2\sin x}-\frac12.
    \]
    Hence \(y=1/2+1/z(x)\) on the component of \(\{x:z(x)\neq0\}\)
    containing zero. The second constant solution \(y=-3/2\) is included
    by \(z=-1/2\).
\end{snippetexample}

\begin{snippetexample}{diffeq-riccati-logarithm-example}{A nonconstant particular solution}
    Consider, for \(x>0\),
    \[
        y'-y^2+2y\ln x-\ln^2x-\frac1x=0,\qquad y(1)=a.
    \]
    The \function \(y_*=\ln x\) is a particular solution because
    \[
        \frac1x-\ln^2x+2\ln^2x-\ln^2x-\frac1x=0.
    \]
    If \(a=0\), the unique solution is \(\ln x\) on \((0,+\infty)\).
    If \(a\neq0\), set \(y=\ln x+1/z\), so \(y'=1/x-z'/z^2\).
    Substitution cancels all logarithmic terms:
    \[
        -\frac{z'}{z^2}-\frac1{z^2}=0,\qquad
        z'=-1,\qquad z=-x+C.
    \]
    The initial value gives \(a=1/(C-1)\), hence
    \[
        y(x)=\ln x+\frac1{-x+1+1/a}.
    \]
    Excluding the pole \(x=1+1/a\) and retaining the component containing \(1\)
    gives the maximal interval
    \[
        I=\begin{cases}
          (0,1+1/a),&a>0,\\
          (1+1/a,+\infty),&a<-1,\\
          (0,+\infty),&-1\leq a<0.
        \end{cases}
    \]
\end{snippetexample}
\end{document}
